\section{Resultados}
\label{sec:resultados}

\ifthenelse{\equal{\resdir}{piloto}}{%
\noindent\fbox{\parbox{0.97\linewidth}{\small\textbf{Atenção:} esta versão mostra o \emph{piloto} reduzido
(CPU; 4 trajetórias de treino por regime, 4 épocas, pré P0, P1, P6 e P9, 3 trajetórias de teste por
regime e 4 blocos por frame). Ele valida o pipeline de ponta a ponta, mas \emph{não} é o experimento
do artigo. Os números finais vêm de \texttt{bash scripts/rodar\_tudo.sh} em GPU, com
\texttt{\textbackslash resdir} trocado para \texttt{exp}.}}\par\medskip}{}

\subsection{Tabelas visuais por instante}
As Figuras~\ref{fig:visual_gliders}--\ref{fig:visual_spots} mostram, para uma trajetória de teste de
cada regime, o instante inicial ($t=0$), dois instantes intermediários ($t=400$ s, na formação do
padrão, em que a variação entre frames é máxima, e $t=5000$ s, no meio da simulação) e o instante
final ($t=10\,000$ s). As linhas mostram o campo $A$ bruto, o campo após o pré-processamento mais
escolhido na condição pré e pós, o GT e, para cada arquitetura, o mapa de erros nas quatro condições
(VP em preto, FP em vermelho, FN em azul). A coluna $t=0$ é a condição inicial e fica fora da
avaliação quantitativa.

\newcommand{\visual}[2]{%
\begin{figure}[p]\centering
\includegraphics[height=0.92\textheight,width=\textwidth,keepaspectratio]{\figres{tabela_visual_#1}}
\caption{Tabela visual do regime #2: entrada, pré, GT e máscaras das três arquiteturas nas quatro
condições (pré/pós escolhidos na validação; legenda \emph{pré / pós}).}
\label{fig:visual_#1}
\end{figure}}
\visual{gliders}{Gliders}
\visual{bubbles}{Bubbles}
\visual{maze}{Maze}
\visual{worms}{Worms}
\visual{spirals}{Spirals}
\visual{spots}{Spots}

\subsection{Tabela de métricas padronizada}
A Tabela~\ref{tab:condicoes123} traz, em ordem sequencial, as condições sem pré e sem pós, somente
pré e somente pós. A condição pré e pós completa está no Apêndice (Tabela~\ref{tab:condicao4}), e a
Tabela~\ref{tab:resumo} resume, por arquitetura, a referência e a melhor configuração de cada
condição. A Figura~\ref{fig:heatmap} mostra as 363 combinações de uma vez: a célula
(nenhum, nenhum) é a referência, a linha ``nenhum'' é somente pós, a coluna ``nenhum'' é somente
pré e o restante é pré e pós.

\tabres{t_resumo}
\tabres{t_condicoes_123}

\begin{figure}[t]\centering
\includegraphics[width=\textwidth]{\figres{heatmap_dice}}
\caption{Dice no teste para cada par (pré, pós) e arquitetura. A cor é a diferença em relação à
referência (sem pré e sem pós, célula destacada).}
\label{fig:heatmap}
\end{figure}

\subsection{Resultados por regime, temporais e de custo}
\tabres{t_regimes}
\tabres{t_temporais}
\tabres{t_custo}

\begin{figure}[t]\centering
\includegraphics[width=\textwidth]{\figres{dice_por_regime}}
\caption{Dice no teste por regime (IC de 95\%) nas quatro condições, por arquitetura.}
\label{fig:regimes}
\end{figure}

\begin{figure}[t]\centering
\includegraphics[width=\textwidth]{\figres{curvas_treino}}
\caption{Perda de treino (contínua) e de validação (tracejada) por época, para cada pré-processamento.}
\label{fig:curvas}
\end{figure}
